%% =====================================================================
%% MATERIAL MOVED FROM PAPER 1 TO PAPER 2
%% Original location: Paper 1, end of Section 5 (Physical realization)
%% Target location: Paper 2, Section 15 (BEC predictions), after the
%%   golden-ratio hyperbolic BEC paragraph
%% Date: 2026-03-19
%% =====================================================================

What would need to be realized is the hyperbolic-plane Green's
function for vortex interactions, $h(r_{\mathbf{H}^2})$, where
$r_{\mathbf{H}^2}$ is the geodesic distance~\citep{Kimura1999}.
In a 2D superfluid hosted on a hyperbolic lattice---for example,
a Bose--Hubbard model implemented in circuit QED---the vortex
interaction would inherit the hyperbolic geometry.  Equivalently,
topological defects in a photonic condensate on a Poincar\'e-disk
resonator array would interact via the hyperbolic Hamiltonian.
Neither realization exists yet.  The prediction $N_{\mathrm{crit}} = 12$
at $R/a = 1$ (and $N_{\mathrm{crit}} \ge 14$ at $R/a = 2$) is
a falsifiable quantitative target for such an experiment.

Three conditions are required for the hyperbolic Hamiltonian
$h = -\ln\sinh(d_{\mathbf{H}^2}/2a)$ to emerge as the effective vortex
interaction.  First, the platform must support a 2D superfluid order
parameter $e^{i\theta}$ on the hyperbolic lattice, so that vortices
appear as topological defects with winding number $\pm 1$ in $\theta$.
Second, in the long-wavelength limit the phase field must satisfy the
hyperbolic Laplace--Beltrami equation $\Delta_g\theta = 0$, whose
Green's function on $\mathbf{H}^2$ is precisely $h$~\citep{Kimura1999};
this follows automatically once the lattice dispersion relation is
hyperbolic.  Third, the curvature radius $a$ and ring radius $R$ must
both be independently controllable: $a$ is set by the lattice geometry
(fixed at fabrication, as demonstrated by~\citealt{Kollar2019}),
while $R$ can be tuned by pinning vortex excitations to engineered
impurities or to the boundary of the resonator array.  In a
Bose--Hubbard model on a Poincar\'e-disk resonator array these three
conditions are simultaneously achievable: the superfluid phase
satisfies the first, the hyperbolic dispersion relation the second,
and current fabrication control over impurity placement is sufficient
for the third.  A realized system of this type would allow direct
measurement of the $N_{\mathrm{crit}}(R/a)$ curve---a signature
inaccessible in flat geometry---by counting the largest stable
vortex ring as $R/a$ is varied.
